\SetPicSubDir{pic/ch2-Review}

\subsection{Literature Review}
\label{sec:LitReview}

In this section, the challenges in current OCC examples are investigated and identified to clarify directions for future research and development.

A systematic review was conducted to summarize existing research on OCC.
\autoref{table:search_criteria} presents the search criteria used to identify relevant articles from the Scopus database. 
The initial search resulted in 309 publications.
The titles and abstracts of these studies were screened to identify relevant studies focusing on thermal comfort control, resulting in 188 papers. 

\begin{table}[pos=!bt]
  \centering
  \input{pic/ch2-Review/criteria}
  \caption{Search criteria for paper selection}
  \label{table:search_criteria}
\end{table}

\section{Static and custom comfort models in operation}
\label{subsec:StaticandCustom}

First, trends in control–option determination methods used in selected OCC studies are analyzed.
Target cases are classified into “temperature threshold–based,” “PMV threshold–based,” and “custom comfort model–based” depending on how control actions are chosen.
The first two types typically rely on a preset indicator—such as a predefined temperature setpoint or a PMV threshold—rather than focusing on occupants’ thermal perceptions. 
In practice, the temperature–threshold approach is a conventional range–based control that does not reflect specific occupant feedback. 
Meanwhile, the PMV–threshold method determines setpoints to bring the environment within a comfortable range based on PMV, representing a collective or typical occupant’s sensation rather than accounting for personal variability. 
Consequently, neither of these two methods truly addresses individual differences among people actually working in the space. 
In contrast, approaches in the “custom comfort model–based” category select control options according to specialized models tailored to specific needs, thereby capturing occupant–specific preferences.

\begin{figure}[pos=!bt]
  \centering
  \includegraphics[width=0.6\linewidth]{\Pic{png}{comfort_determining_all}}
  \caption{OCC studies by control setpoint determination methods (All)}
  \label{fig:metrics_all}
\end{figure}

\begin{figure}[pos=!bt]
  \centering
  \begin{minipage}{0.49\textwidth}
    \centering
    \includegraphics[width=\linewidth]{\Pic{png}{comfort_determining_Field}}
    \caption{Field Studies}
    \label{fig:metrics_field}
  \end{minipage}\hfill
  \begin{minipage}{0.49\textwidth}
    \centering
    \includegraphics[width=\linewidth]{\Pic{png}{comfort_determining_simu}}
    \caption{Simulation}
    \label{fig:metrics_simu}
  \end{minipage}
\end{figure}

\autoref{fig:metrics_all}–\autoref{fig:metrics_simu} show an increasing trend in simulation cases, whereas the increase in real–building cases is less pronounced.
Under “thermal comfort” and “OCC”–related keywords, many examples still use temperature or PMV threshold control, especially in simulations.
However, field studies utilizing custom comfort models remain limited.

\section{Occupancy conditions in existing OCC studies}

\subsection{Target space occupancy count}

This and the following subsection investigate occupancy–related contexts focusing on OCC examples that correspond to “custom comfort model–based” in \autoref{subsec:StaticandCustom}.
Because the number of occupants in a space (i.e., occupancy scale) can profoundly impact control strategies, this subsection explores relationships to experimental settings, control resolution, and algorithms.

\begin{figure}[pos=!bt]
  \centering
  \includegraphics[width=0.7\linewidth]{\Pic{png}{Occu_histo_FieldSimu}}
  \caption{Occupancy count of target spaces in field and simulation settings}
  \label{fig:Occu-Field}
\end{figure}

\autoref{fig:Occu-Field} shows the ratio of field and simulation studies along the occupancy scale.
For single–occupant spaces, both simulation and field studies are common.
However, for multi–occupant spaces, the ratio of field studies is relatively higher—relevant for real operations (e.g., offices), highlighting the need for OCC strategies that address shared usage dynamics.

\subsection{Control resolution}

The relationship between occupancy scale and control resolution (HVAC control granularity) is examined, and strategies are classified into three categories: room control, group control, and personal control.

In room control, the system determines options (e.g., VAV setpoint) for the entire room.  
For single–occupant rooms, applying comfort models is straightforward.
Multi–occupant rooms require careful consideration of diverse preferences.

In group control, the system conditions room subzones via additional devices (e.g., ceiling fans or radiant panels), enabling more flexible satisfaction of target groups based on merged comfort models \cite{ono_effects_2022}.

In personal control, the system adjusts to individual preferences, often with personal devices \cite{andre_user-centered_2020,xu_towards_2022,aryal_intelligent_2021}.
This review also includes seat/room assignment approaches \cite{nagarathinam_user_2021,sood_spacematch_2020} that match occupants to environments.

\begin{figure}[pos=!bt]
  \centering
  \includegraphics[width=0.7\linewidth]{\Pic{png}{ControlReso_Occu}}
  \caption{Control–resolution trends across occupancy levels}
  \label{fig:Occu-conreso}
\end{figure}

\autoref{fig:Occu-conreso} indicates most studies adopt room–level control—even for multi–occupant rooms—likely because VAV/FCU/DX room–level systems are common in practice.
Fewer studies consider group/personal control using additional equipment (e.g., ceiling fans \cite{lei_practical_2022}, personal devices \cite{boudier_analysis_2022,aryal_intelligent_2021}, radiant panels \cite{lee_implementation_2019}).

\subsection{Occupant-centric control in dynamic occupancy context}

In dynamic occupancy environments, occupancy density and distribution fluctuate throughout the day, making static strategies ineffective \cite{ono_occupant-centric_2024,sood_spacematch_2020}.
A strategy optimized for one scenario becomes suboptimal for another, leading to discomfort, reduced productivity, and higher energy use.

% \textcolor{red}{[Add more ABW studies]}

Two control directions under dynamic seating emerge: (i) space–matching and (ii) dynamic–occupancy–based control.

\subsubsection{Space–matching method}

The idea is to provide a range of indoor environments rather than a single uniform condition, then assign occupants to maximize preference–environment matches.
Assignments can be at the room level or seat level.

Room–level examples include \cite{sood_spacematch_2020,ding_enhancing_2023,nagarathinam_user_2021,huang_space-matching_2025,huang_improving_2023}.
Ding et al.~(2023) \cite{ding_enhancing_2023} simulated comfort/energy impacts of room–selection strategies with different setpoints.
Ono et al.~(2024) \cite{ono_occupant-centric_2024} extended objective dimensions (preference, availability, function type, travel distance, whole–building energy).

Some studies tested in real buildings: Huang et al.~(2023) \cite{huang_improving_2023} examined satisfaction in meeting rooms; Sood et al.~(2020) \cite{sood_spacematch_2020} clustered preferences and suggested seats accordingly.

Seat–level examples include \cite{nagarathinam_user_2021}; similarly, \cite{huang_space-matching_2025} combined CFD with seating policies and found dynamic seating reduces energy in simulations.

While space–matching can incorporate diverse factors (lighting, visual/acoustic conditions, equipment availability), weights differ across workers, and selections may deviate from system suggestions.

\subsubsection{Dynamic–occupancy–based control}

As noted in reviews \cite{xie_review_2020}, integrating diverse preferences across multiple occupants remains underexplored.
This section focuses on how comfort models are incorporated into control.

In some studies on the effectiveness of OCC, such as \cite{jung_energy_2020}, occupancy is often assumed to be static.  
However, in today’s dynamic working environments—where hybrid working, combining remote and office-based work, has become commonplace\cite{redmond_rise_nodate} —this static assumption may underestimate the impact of occupancy characteristics.

\autoref{table:OccuIntegration} compares field studies that integrate multiple personal comfort models within the same control zone.

The “Occupancy Adaptation” column indicates each framework’s ability to handle different occupancy patterns; “dynamic” adjusts to real–time combinations, while “static” assumes fixed configurations. 
Only \cite{lei_practical_2022} considers transitions among combinations by pre–training multiple RL models for frequent patterns—data–and compute–intensive.

\begin{table}[pos=!bt]
  \centering
  \begin{tabular}{|l|c|c|c|c|c|c|}
\hline
\textbf{Paper} & \textbf{Year} & \textbf{Resolution} & \textbf{Per unit} & \textbf{Algorithm} & \textbf{Occupancy adaption} \\
\hline
\cite{lei_practical_2022} & 2022 & Group & 4 people& RL & Dynamic  \\
\cite{lee_implementation_2019} & 2019 & Group & 2 people & MPC & Static \\
\cite{li_indoor_2019} & 2019 & Room & 4 people & MPC & Static \\
\cite{li_personalized_2017} & 2017 & Room & 7people & RB & Static \\
\cite{li_indoor_2023} & 2023 & Room & 4 people & RL & Static \\
\cite{aguilera_thermal_2019} & 2019 & Room & 16 people & RB & Static \\
\cite{li_experimental_2021} & 2021 & Room & 6 people & RB & Static \\
\hline
\end{tabular}
  \caption{Field OCC studies integrating multiple personal comfort models}
  \label{table:OccuIntegration}
\end{table}

\begin{figure}[pos=!bt]
  \centering
  \includegraphics[width=\linewidth]{\Pic{png}{DynamicOccu}}
  \caption{Schematic of dynamic occupancy}
  \label{fig:Upd-Field}
\end{figure}

Historically, tracking personal–level occupancy was difficult in typical BMS; this has changed with IoT systems enabling occupancy/location tracking \cite{soleimanijavid_challenges_2024}.
Thus, \autoref{table:OccuIntegration} underlines the potential of real–time occupancy for OCC frameworks as an underexplored domain.

\section{Comfort models in occupant-centric control}

% \subsection{Developing group/zone comfort models}

% As suggested by \autoref{fig:Occu-conreso}, interest is growing in group/zone comfort models comprising multiple occupants within the same control area.
% Two approaches exist:
% \begin{enumerate}
%   \item Train individual PCMs for multiple occupants and aggregate them.
%   \item Train group models without modeling individual differences.
% \end{enumerate}

% \textcolor{red}{[Add analysis on group comfort vs. aggregation of PCMs]}

\subsection{Updating comfort models during operation}
\label{sec:UpdComMod}

For long–term OCC, updating comfort models becomes inevitable.
Models calibrated for specific seasons/conditions may degrade as clothing insulation and metabolic rates change.
As noted by \cite{aguilera_thermal_2019}, using fixed PCMs over long periods without updates can reduce satisfaction, particularly when initial data are insufficient.

\begin{figure}[pos=!bt]
  \centering
  \includegraphics[width=0.7\linewidth]{\Pic{png}{Update_fldsimu}}
  \caption{Ratio of studies updating comfort models by experiment setting}
  \label{fig:Upd-Field}
\end{figure}

\begin{figure}[pos=!bt]
  \centering
  \includegraphics[width=0.75\linewidth]{\Pic{png}{ConAlgo_Upd}}
  \caption{Ratio of studies updating models by control algorithm}
  \label{fig:Upd-Algo}
\end{figure}

Notably, even though reinforcement learning (RL) updates action values with state transitions, the comfort–model component is seldom revised; \cite{lei_practical_2022} is a rare exception.
Most RL work focuses on learning HVAC control policies and building dynamics, not updating the comfort models themselves.

\begin{table}[pos=!bt]
  \centering
  \input{pic/ch2-Review/UpdateStudies}
  \caption{Studies updating comfort models during operation}
  \label{table:UpdateStudies}
\end{table}

\subsection{State-of-the-art comfort model development methods}

Recent deep learning advances enable more efficient and accurate models, reducing survey effort.

Transfer learning leverages knowledge from a source dataset to efficiently train new models: pretrain on a large source dataset, then adapt to a smaller target dataset.  
For instance, \cite{gao_transfer_2021,somu_hybrid_2021} used the ASHRAE Global Thermal Comfort Database \cite{foldvary_licina_development_2018} and improved accuracy compared to training on limited data.

Combining transfer learning with active learning can further reduce survey counts by querying the most informative samples \cite{tekler_data-efficient_2024}.
\cite{horikoshi_developing_2024} examined how the quantity of survey data affects comfort and energy across different control settings.

Despite these advantages, most applications target individual PCMs.  
In principle, OCC does not always require finely tuned PCMs for every occupant: in real buildings with many users, zone- or group-level models may provide sufficient accuracy while further reducing the survey burden.
This is because the marginal value of distinguishing each individual diminishes when the occupant pool is large and heterogeneous\cite{jung_energy_2020}.
Conversely, under states such as small groups or high turnover, the added resolution of individual PCMs can be highly informative and justify additional survey effort.
Therefore, while TL and AL are primarily proposed as methods to ease the creation of personal models, a practical OCC framework must balance between individual- and group-level modeling, depending on occupancy context.

\section{Summary of literature review and identified knowledge gap}

Key gaps observed in OCC research:
\begin{itemize}
  \setlength{\itemsep}{0pt}\setlength{\parskip}{0pt}
  \item Single–occupant scenarios predominate in both field and simulation studies. 
  \item Even with multiple occupants, room–level control is the default.
\end{itemize}

Operational adaptability gaps:
\begin{itemize}
  \setlength{\itemsep}{0pt}\setlength{\parskip}{0pt}
  \item Few studies update comfort models during operation. 
  \item Real–time occupancy is rarely incorporated within simple, practical frameworks.
  \item No framework effectively integrates individual PCMs in real time under continuous occupancy fluctuations.
\end{itemize}

On comfort–model development:
\begin{itemize}
  \setlength{\itemsep}{0pt}\setlength{\parskip}{0pt}
  \item Transfer/active learning reduce workload for individual PCMs. 
  \item For group models, aggregating many PCMs may be impractical; scalable alternatives are needed for real–world deployment.
\end{itemize}

\label{sec:researchaim}
As outlined in the introduction, this research addresses the critical gaps limiting the practical application of Occupant-Centric Control.

\subsection{Research Problem and aim}
There is no clear guidance on when personal comfort models are preferable at the room-control resolution, especially under conditions with dynamic seating and turnover.  
Specifically, this study aims to Provide principled guidance for selecting room-level comfort model strategies under real-world, dynamic occupancy, and

\subsection{Research Questions and Hypotheses}
The investigation is guided by specific research questions (RQs) and hypotheses (Hs), which are structured around the two-step experimental approach.
This experiment examines how occupancy patterns affect the performance gap between dynamic and static comfort models.

\textbf{Hypothesis 1 (H1):}  
The effectiveness of static GCMs diminishes as occupancy conditions become unstable.

\textbf{Research Question 1 (RQ1):}  
To what extent do occupancy patterns (average occupancy, variability, and user turnover rate) influence the relative effectiveness of PCMs over static GCMs in office environments?

\textbf{Hypothesis 2 (H2):}  
Leveraging real-time occupancy with PCMs improves thermal comfort and energy efficiency compared with static, group-based control at the room level.

\textbf{Research Question 2 (RQ2):}  
To what extent does using real-time occupancy and a dynamic PCM-based framework improve room-level OCC in terms of comfort and energy, relative to a static, group-based approach?